% Pursuing Stacks (cotes 134-2 to 134-8): archive cote 134-5, site page 125 (PDF page 126), typescript page 312.
% typescript. Draft compared: ps5.tex:987-1065.
% AI-transcribed by Claude Opus 5.5 (claude-opus-5-5, adaptive thinking, effort high), Message Batches API, batch ps3 (msgbatch_01XdFmuVVBAcgc7EFxdtCqsj), 2026-11-12; reviewed against both supplied scan halves; see CHECK_DECISIONS.tsv.
% Hand-drawn arrows are marked blue; typed objects remain plain.
\scanpage{134-5/125}
\mspage{312}

following useful statement,  which is lacking in SGA IV Chap.I
(compare loc. cit.  prop.7.2 , page 47,  giving part of the story):
\hand{\uline{\textcolor{black}{Proposition}}}  Let M be a U-category,  A a small \inserted[typed]{full} subcategory,  i: A $\hand{\to}$ M
the inclusion functor,  hence a functor
\[ i^{*} : M \hand{\longrightarrow} A^{\wedge} \quad . \]
For any object X of M,  we consider the family \xout{\unsure{$\Phi$}} $F_X$  of all arrows
in M with target X,  source in A.  The following conditions are equiva-%
lent:

(i)  For any X in M,  the family $F_X$ is strictly epimorphic.

(ii) For any X in M,  the family $F_X$ is universally strictly epimor-%
phic.

(iii)  For any X in M ,  $F_X$ \xout{induces an isomorphism}\inserted[typed]{turns X into a direct}  limit of the
composition functor  $A_{/X} \hand{\to} A \hand{\to} M$ ,  i.e.
\[ X \hand{\xleftarrow{\sim}} \underset{A/X}{\underset{\hand{\longrightarrow}}{\mathrm{lim}}} a \qquad . \]
(iv)  The functor $i^{*}$ is fully faithful.

Proof left to the reader (who may consult loc. cit. for \xout{(i)}
(i)$\hand{\Longrightarrow}$(iii)$\hand{\Longleftrightarrow}$(iv) ,  so that only (i)$\hand{\Rightarrow}$(ii) is left to prove).

\hand{\uline{\textcolor{black}{Definition}}}  When the equivalent conditions above are satisfied, we'll
say that A is a \hand{\uline{\textcolor{black}{strictly generating}}} subcategory of M.

NB. The notion makes a sense too without assuming  A to be small,
nor M to be a U-category,  by passing to a larger universe  (it is imme-%
on (i) or (iii)  that the condition does not depend on the choice of the
universe\hand{)}.

\hand{\uline{\textcolor{black}{\over[typed]{O}{C}orollary}}}  If A is strictly generating in M,  then so is any larger
full subcategory B.

This is clear by criterion (ii)  (whereas it is'nt by any one of
\margin{\tikz[x=.45em,y=.55em]{\draw[->,color=blue] (1.6,0)--(0,0)--(0,2.2);}}the other three criteria !).

Digressions are ended,  let's resume our review!

\xout{\unsure{91}}
\margin{\tikz[x=.45em,y=.55em]{\draw[->,color=blue] (1.8,2)--(0,2)--(0,0);}}91.   I see three more questions to review -  presumably \deleted{it}\inserted{they} will be
a lot shorter than the last !

7)  \hand{\uline{\textcolor{black}{Description of homotopy types}}} ``\hand{\uline{\textcolor{black}{of finite type}}}''\hand{,} in terms of an
elementary modelizer $A^{\wedge}$.  In\hand{\tikz[x=.22em,y=.55em,baseline=-.2ex]{\draw (0,1.4)--(1,-.5);}}terms of the modelizer (Spaces), a natural
finiteness condition on a homotopy type is that it may be described (up

\margin{125}
